\section{Results on SummEval}
\label{sec:summeval}

All numbers in this section are computed on the 50 held-out SummEval articles (800 summaries). Table~\ref{tab:confound} is the main result: for every metric it gives the cost, the raw and the partialled mean correlation, the correlation with the copy rate, the partialled correlation per dimension, and whether the metric is on the cost--quality Pareto frontier of either axis. Appendix~\ref{app:tables} gives the raw per-dimension correlations with bootstrap intervals and the system-level correlations.

\subsection{The extractive shortcut}

The last row of Table~\ref{tab:confound} scores the copy rate as if it were a metric. Its mean raw correlation of .271 is higher than that of BERTScore (.269), EmbedScorer, SMART-Model, MoverScore, BARTScore and all lexical metrics: ten of the fourteen baselines are beaten by a ratio of shared bigrams. The human ratings are themselves correlated with the copy rate (Spearman .21 for coherence, .38 for consistency, .31 for fluency and .20 for relevance over all of SummEval), and every metric that looks at the source inherits part of that correlation. The lead-window controls, which match the summary against the first three or five source sentences with ROUGE-L, are the clearest case. Lead-5 reaches a raw $\bar\rho$ of .292, higher than BERTScore, GPTScore and AlignScore, but its correlation with the copy rate is .708 and its partialled $\bar\rho$ falls to .147.

Controlling for the shortcut reorders the pool. Metrics that see only the reference, such as BERTScore ($\rho_{\mathrm{copy}}=.027$), lose nothing, because copying from the source cannot leak into them. Source-based similarity metrics lose most: LGS drops from .299 to .167, CCM from .248 to .142, and the lead controls by half. The learned and model-based evaluators lose least: UniEval drops from .473 to .424 and G-Eval from .327 to .291. The partial axis also separates metrics that the raw axis puts close together: LGS and Lead-5 have raw correlations close to those of GPTScore and AlignScore (.299 and .292 against .290 and .282), but lower partial ones (.167 and .147 against .194 and .200).

\subsection{LNC against the pool}

LNC reaches a raw $\bar\rho$ of .433 and a partial $\bar\rho$ of .350 (95\% interval [.313, .385]). It is second to UniEval on both axes and ahead of every other metric, including the local G-Eval judge (.327 raw, .291 partial). Its loss from raw to partial (.083) is larger than UniEval's (.049) and G-Eval's (.036), and its correlation with the copy rate (.550) is high: LNC still uses copying, but most of its agreement with the human ratings survives when copying is controlled for.

Per dimension, LNC is balanced where its competitors are not. Its partial correlations are .350, .370, .340 and .339 for coherence, consistency, fluency and relevance. UniEval is much stronger on coherence (.550) and weaker on consistency (.289). G-Eval is strong on consistency (.408) and weak on fluency (.149). BERTScore is strong on coherence and relevance (.425, .420) and close to zero on consistency and fluency. AlignScore and NLIG are consistency metrics: .318 and .413 on consistency, at most .295 elsewhere.

Table~\ref{tab:lnc_comparisons_partial} tests these differences on the partial axis with paired bootstrap comparisons over the 36 cells of LNC against nine baselines. LNC is significantly better in 22 cells, tied in 12 and worse in 2, both against UniEval: on coherence ($-.200$) and fluency ($-.086$). Against UniEval it is better on consistency ($+.082$), and the difference on relevance is not significant. Against G-Eval it is better on fluency ($+.190$) and tied on the other three dimensions. Against AlignScore and NLIG it ties on consistency, their own dimension, and is better on coherence and relevance. It is better than LGS, Lead-5 and BARTScore on all four dimensions. On the raw axis the picture is the same, with 25 wins, 10 ties and one loss, to UniEval on coherence (the table is rendered by the released code).

\subsection{Rank tests over the whole pool}

The Friedman test (Table~\ref{tab:friedman}) treats each article as a block and ranks the metrics by their within-article correlation with the human ratings. The null hypothesis of equal average ranks is rejected on every dimension ($\chi^2_F$ between 137.6 and 314.1 with 18 degrees of freedom, $p<.001$). LNC and G-Eval have the two highest average ranks on consistency, and LNC has the second highest on coherence and fluency, after UniEval. Holm-corrected Wilcoxon tests separate LNC from UniEval on coherence and consistency, where UniEval ranks higher on coherence and lower on consistency, and from G-Eval on fluency, but not from G-Eval on coherence, consistency or relevance, nor from UniEval on fluency and relevance. Because this test uses within-article correlations, it is insensitive to differences between articles; that it agrees with the pooled analysis shows that LNC's advantage is not an artefact of pooling.

\subsection{Cost}

LNC needs 216\,ms per summary: 106\,ms for the language model (one source pass per article and five candidate variants), 85\,ms for NLI and 25\,ms for the per-token pass of the splice features. That is less than UniEval with all four dimensions (343\,ms), about the same as BERTScore (214\,ms), and 82 times less than the local G-Eval judge (17.8\,s for four dimensions with three samples each). LNC is on the Pareto frontier of both axes in Table~\ref{tab:confound}: no cheaper metric correlates better, raw or partialled. On the partial axis the frontier runs from the lead controls and METEOR through GPTScore, AlignScore and BERTScore to LNC and UniEval; on the raw axis it consists of the lead controls, LGS, LNC and UniEval. The comparison with UniEval needs one qualification: UniEval scores relevance against the reference, which LNC does not have. Among reference-free metrics LNC is the most accurate at any cost below that of an LLM judge, and more accurate than the LLM judge we could run.
